Use $A=(12,0),B=(6,6\sqrt3),C=(-6,6\sqrt3),F=(6,-6\sqrt3)$. Then $P=(6,-3\sqrt3)$. Parametrising $CP$ and $AF$ and equating coordinates yields $Q = (\dfrac{54}{7} , - \dfrac{30 \sqrt{3}}{7})$. Thus $Q - P = (\dfrac{12}{7} , - \dfrac{9 \sqrt{3}}{7})$ and $\boxed{P Q = \dfrac{3 \sqrt{43}}{7}}$.
\begin{asy}
import graph;
usepackage("amsmath");
import geometry;
size(7cm);
pair[] p={(12,0),(6,6sqrt(3)),(-6,6sqrt(3)),(-12,0),(-6,-6sqrt(3)),(6,-6sqrt(3))}; pair labelCentre=(0,0);for(pair labelVertex:p)labelCentre+=labelVertex;labelCentre/=p.length; draw(p[0]--p[1]--p[2]--p[3]--p[4]--p[5]--cycle); dot(p[0]);label("$A$",p[0],unit(p[0]-labelCentre));dot(p[1]);label("$B$",p[1],unit(p[1]-labelCentre));dot(p[2]);label("$C$",p[2],unit(p[2]-labelCentre));dot(p[3]);label("$D$",p[3],unit(p[3]-labelCentre));dot(p[4]);label("$E$",p[4],unit(p[4]-labelCentre));dot(p[5]);label("$F$",p[5],unit(p[5]-labelCentre));pair P=(6,-3sqrt(3)),Q=(54/7,-30sqrt(3)/7);draw(p[1]--p[5]);draw(p[2]--Q);label("$P$",P,(1,1));label("$Q$",Q,(1,-1));
\end{asy}

